\section{Introduction}

The relationship between media and government is a critical component of electoral accountability \citep{Audits-Ferraz-QJE, PolInfo-Enriquez-JEEA, Marshall}. A well-functioning media holds significant power to hold government actions accountable and inform the public \citep{Stromberg-MediaCoverage}. Crucially, incumbent governments are fully aware of this impact, prompting them to devote substantial resources to control the information environment \citep{Autocracy-Guriev-JPE}. In democracies where direct coercion is costly, governments often resort to subtler forms of influence, such as allocating public resources to reward friendly outlets and discipline critics \citep{ArgCorr-DiTella-AEJAE, MediaCaptures-Zeidl-Econometrica}.

However, persuasion is not monotonic in control. As established by the literature on Bayesian persuasion \citep{MediaBias-Gentzkow-JPE, PersuasionEmpirical-Dellavigna-AR}, rational audiences discount information from sources perceived as biased. This creates a dilemma for the incumbent: total censorship ensures a favorable message but destroys the informational value of the signal, leading to citizen disengagement. To persuade effectively, the government may need to commit to a signal structure that allows for occasional negative realizations to make favorable coverage credible.

This paper studies this \textit{credibility--control trade-off} in the context of Mexico's daily presidential press conference, \textit{La Mañanera}. Following his 2018 victory, President Andrés Manuel López Obrador (AMLO) established this direct-to-public communication channel to set the daily agenda. Initially, access was effectively discretionary, and the equilibrium featured a concentration of friendly questions. We document a decline in digital engagement in the run-up to 2022, consistent with audiences treating the conference as less informative. In a stark shift, the administration introduced a seat lottery in April 2022, randomizing front-row access within media strata and thereby introducing an exogenous source of variation in who is most likely to be called to speak.

I propose that this institutional change functions as a commitment device in a persuasion game. By delegating part of access to a randomization mechanism, the government accepts the risk of on-air dissent---especially in unfavorable states of the world, when market incentives for criticism are high---in exchange for restoring the credibility of the conference.

A key feature of this setting, which distinguishes it from the seminal work of \citep{PersuasionEmpirical-Dellavigna-AR}, is that the government shapes the information environment through two main instruments: access to the conference and the allocation of advertising. Access determines which voices become publicly visible, while advertising affects the incentives media outlets face when deciding whether to ask critical questions. Media outlets, which differ in editorial stance, audience size, and market incentives, weigh the benefits of visible dissent against the potential costs of losing future government support. This framework helps explain how the government can preserve influence over the public narrative without fully eliminating criticism. It also provides a clear bridge to the empirical strategy by identifying the key relationships between access, dissent, market rewards, and advertising discipline.

To test these implications empirically, I use rich administrative, textual, and social media data. The analysis combines transcripts from 1,423 conferences, administrative records of government advertising spending from COMSOC, Google Trends data as a proxy for audience demand, and daily presidential approval ratings from Mitofsky's AMLO Tracking Poll. The identification strategy exploits the random variation induced by the seat lottery: seat location is randomized within media strata and strongly predicts the probability of being called, providing an exogenous source of variation in on-air exposure. I use this variation to identify the market reward for criticism, the implied punishment schedule, and the effects of visible dissent on persuasion and public attention.

The rest of the research proposal is organized as follows. Section \ref{sec:lit_review} provides a detailed description of the literature related to this topic and the historical and political Mexican context. Section \ref{sec:model} provides a theoretical model. Section \ref{sec:data} outlines a description of all the data used in the analysis, as well as an explanation on how I build each variable. Section \ref{sec:emp_strategy} presents the empirical strategy. Section \ref{sec:results} presents preliminary results and descriptive statistics. Section \ref{sec:conclusion} concludes.
